% !PRINTED 302
\weberpage{302}

%<<CONT
there would also be one among them which, without $r, s$ vanishing, becomes identical with $V$. But this requires

\begin{align*}
\alpha\equiv\pm\alpha g^{r+s}, &\qquad \beta\equiv\pm\beta g^{r-s}\\
\gamma\equiv\pm\gamma g^{-r+s}, &\qquad \delta\equiv\pm\delta g^{-r-s}\quad(\mathrm{mod}\ p),
\end{align*}

where in all four formulae the upper or the lower signs hold, and hence, since neither $\alpha$ and $\delta$ nor $\beta$ and $\gamma$ vanish simultaneously,

\[ r+s\equiv 0,\qquad r-s\equiv 0\qquad(\mathrm{mod}\ p-1), \]
or
\[ r+s\equiv\frac{p-1}{2},\qquad r-s\equiv\frac{p-1}{2}\qquad(\mathrm{mod}\ p-1), \]

and so in both cases

\[ r\equiv 0,\qquad s\equiv 0\ \left(\mathrm{mod}\ \frac{p-1}{2}\right),\qquad\text{q. e. d.} \]

In the form (16) there are therefore contained $\left(\dfrac{p-1}{2}\right)^2$ \gesperrt{different substitutions}. If the group $\mathfrak{K}$ is not yet exhausted by them, one chooses a substitution $V'$ not contained in (13), (14), (16) and forms in the same way the sequence

\weq{17}{C^{r}V'C^{s},}

whose substitutions are different both among themselves and from those contained in (13), (14), (16); and one proceeds in this way until the whole group $\mathfrak{K}$ is exhausted.

If we denote by $q$ the number of the sequences (16), (17), $\dots$ so formed, then the degree of the group $\mathfrak{K}$ comes out:

1. if $B$ does not occur in $\mathfrak{K}$:

\[ =\frac{p-1}{2}+q\left(\frac{p-1}{2}\right)^2; \]

2. if $B$ occurs in $\mathfrak{K}$:

\[ =p-1+q\left(\frac{p-1}{2}\right)^2, \]

and by the requirement of our problem this number is to be

\[ =\frac{(p-1)(p+1)}{2} \]

From this, however, it follows that case 1. is impossible, for in this case we should necessarily have

\[ q\left(\frac{p-1}{2}\right)=p, \]
%CONT>>
